\subsection{An explicit renderer for private images and intents}
\label{sec:opaque-rendering}

The alias lemma preserves a renderer that can read every original payload.
A component reduction additionally needs to execute record management when
the selected KEM challenger supplies public outputs and a shared-secret
challenge without supplying the selected random tape. We construct that
record-management algorithm for a finite image/intent interface. Its tables
contain immutable handles, known fields and locally sampled sealing coins.
The protocol's public computations are supplied by a common functional
interface $F$, instantiated by the original-operation adapter. This section
holds $F$ fixed and isolates the image/intent contribution to the record
connection in Proposition~\ref{prop:native-local-accounting}.

\paragraph{Typed record interface.}
An opaque block is an immutable fixed-length private value, represented by a
handle and its known length. Blocks may be correlated and may themselves be
retained sealed reservations. A private image recipe contains at least one
such block; an atom-free recipe is fully known. Local image-encryption and
intent-authentication keys are available to the simulator. Fix a public
budget $M$ for all fresh image-sealing invocations, including failed attempts
and atom-free images. Exact copies and exact retries retain their invocation.
The following operations specify the interface.

\begin{enumerate}
  \item \emph{Seal.} Before sampling its nonce, invocation $i$ fixes the
  recipe, its length and the non-nonce header fields. It samples independent
  $N_i\gets\{0,1\}^{n}$ and privately constructs
  \[
    B_i=\mathsf{Hdr}_i\Vert N_i\Vert C_i,
  \]
  where $C_i$ is the specified deterministic local authenticated encryption
  under that nonce. Local encryption and intent authentication use their
  prescribed fixed algorithms without additional shared-SHA3 calls. All
  admitted image formats place the nonce at the same offset. The reply
  consists of the image handle, length and declared non-nonce metadata;
  $N_i$, $C_i$ and the recipe's private blocks remain private.
  \item \emph{Intent.} A complete canonical prefix $P$ fixes the tag,
  binding, predecessor and the image-head value used in this execution.
  Its tag determines its length. The intent bytes are
  \[
    W_{i,P}=P\Vert B_i\Vert
       \mathsf{MAC}_{k_i}(P\Vert B_i),
  \]
  with one fixed authentication key $k_i$ for image $i$. An exact repeat
  retains the same $(i,P)$ entry. Image and intent tags are disjoint.
  \item \emph{Public commitments.} For a designated private image, the
  ideal-commitment experiment returns independent uniform 256-bit values
  cached by $i$ for its head and by $(d,i,P)$ for its intent digests. The
  domain $d$ ranges over the five admitted intent-digest domains. Atom-free
  objects retain their real hashes and ordinary raw-byte outputs. Later
  validation and zero-value checks use the actual replies in that execution.
  \item \emph{Equality and storage.} Equality queries compare complete
  retained image or intent handles. Copy, retain, persist, delete, crash
  and resume update an explicit durable-state directory. Storage choices,
  errors and branch decisions depend on available metadata, lengths,
  commitments and $F$ outputs. A failure after commit preserves that commit.
  \item \emph{Acquisition.} Atom-free objects are rendered from their known
  inputs. For an opaque object, the admitted successful acquisition emits
  a computable event that permanently falsifies the fixed-candidate
  eligibility before its raw bytes are returned. The stopped score is zero,
  as in Corollary~\ref{cor:target-acquisition-stop}. Unavailable-object requests
  retain their ordinary error result.
\end{enumerate}

All other public computation uses these replies, known inputs, $F$, and the
shared public oracles. This is a classical record interface; the adversary
may have quantum private state and quantum access to those public oracles.
The successful-acquisition condition is supplied by the recovery experiment's
existing event rule. For a fixed target session $s$ and cut $c$, acquisition
of an original version $v$ with $\mathsf{subject}(v)=s$ and
$c<\mathsf{save}(v)$ supplies such an event. Its original-record provenance
determines that save time, including when a later aggregate carries the same
version. Lemma~\ref{lem:adaptive-origin-stop} gives a provenance certificate
valid before the final cut is chosen. A separately permitted projection is
described by its own output operation.

\paragraph{Table algorithm.}
The simulator stores an image table indexed by fresh sealing invocation,
an intent table indexed by image and complete prefix, a commitment cache,
and the directory with its provenance and acquisition ledgers. Each new
image entry records the recipe handles, length, known fields and a sampled
nonce. Each new intent entry records $(i,P)$ and its known length. The
simulator samples designated-private commitments on first use and retrieves
them on exact retry; it constructs and hashes atom-free objects normally.
Equality is structural: images compare equal exactly when their indices
match, and intents compare equal exactly when their pairs $(i,P)$ match.
The storage directory determines the state recovered after a crash, including
an acknowledged or uncertain write that actually committed. At the stipulated
permanent-failure acquisition event the computation returns zero. Every other
transition executes the common public code and $F$. Thus the algorithm uses
no selected opaque block bytes.

\begin{lemma}[Structural equality of complete records]
\label{lem:record-structural-equality}
If the nonces of fresh image invocations are pairwise distinct, then
\[
 B_i=B_j\iff i=j,\qquad
 W_{i,P}=W_{j,P'}\iff(i,P)=(j,P').
\]
\end{lemma}
\begin{proof}
One immutable image handle retains one byte string. Two different image
invocations have different nonces at the same explicit offset, so their
complete wires differ. Intent tags first separate different formats. For
one tag the prefix length is fixed: unequal prefixes differ before the
target image, while equal prefixes with different images differ at the
target's nonce. The same image and prefix give the same deterministic MAC
under the fixed key. Disjoint tags also separate images from intents.
This argument permits coincident private block values and ciphertext
components; the explicit nonces distinguish the complete images.
\end{proof}

Let $I$ be the materialized ideal-commitment experiment, which retains real
private bytes and answers equality by comparing them. Let $T$ be the table
algorithm above. Both use the stipulated fixed-candidate stopping rule and
the same public-oracle model and interface $F$.

\begin{theorem}[Opaque image/intent simulation]
\label{thm:opaque-rendering}
For at most $M$ fresh seals with independent $n$-bit nonces, define
\[
 \delta_{\rm nonce}=\min\{1,\tbinom M2\,2^{-n}\}.
\]
For every efficient terminal bit $D$ computed from the public view, available
metadata and event ledger, $F$ outputs and its own independent test coins
(including its retained application challenge bit),
\[
  |\Pr[I:D=1]-\Pr[T:D=1]|\leq\delta_{\rm nonce}.
\]
For a score $Z\in[-1,1]$ computed from the same available inputs,
$|\mathbb E_I Z-\mathbb E_T Z|\leq2\delta_{\rm nonce}$.
\end{theorem}
\begin{proof}
Couple the fresh sealing nonces, private-commitment replies, scheduler choices,
the state of $F$, and the public-oracle transformations. Before the first
nonce collision, maintain agreement of the directories, known fields,
immutable references, public bytes, statuses and semantic ledgers.
Sealing adds the same metadata and nonce to corresponding entries; its
private ciphertext is confined to the materialized entry. Intent creation
adds the same complete prefix and target reference. A fresh private
commitment gets the coupled reply, and a repeat consults the same cache.
Atom-free commitments and acquisitions use identical materialized inputs.
Lemma~\ref{lem:record-structural-equality} gives agreement of equality replies.
The common public computation consequently makes the same transitions and
oracle calls. Copying, committed or failed writes, deletion, crash and
resume preserve the directory invariant. An opaque acquisition emits the
same permanent-failure event before either computation needs to return its
unknown bytes, and both stopped scores are zero.

The next nonce is independent uniform even when the decision to seal is
adaptive. Across at most $M$ invocations, a union bound charges
$\binom M2 2^{-n}$ for any collision. Coupled terminal bits differ only on
that event; bounded scores can differ by at most two. Identical classical
replies and shared-oracle transformations also preserve a quantum
adversary's state at every coupled good prefix. The argument uses no
measurement of a quantum query.
\end{proof}

\paragraph{Connecting real commitments.}
After replacing complete-wire equality by structural equality, let
$\epsilon_{\rm com}$ bound the bit-distinguishing gap between real private
record digests and the independently cached commitments of the typed
interface. Both games run the same public code, ordinary known-record
computations, available metadata and declared stopping rule. The comparison
counts the full distinguishing computation against their shared public oracle.
It is a cryptographic input to the following composition, separate from the
explicit table simulation. Appendix~\ref{app:primitive-inputs} records the
concrete multi-record bound under fixed-input injectivity, per-nonce
confinement and the prescribed local encryption/authentication algorithms.

\begin{corollary}[Real-record to table-algorithm bound]
\label{cor:real-to-opaque}
For the declared fixed-candidate interface and the commitment conditions
above, let $R$ be the real-digest materialized experiment. With $n=192$,
\[
 |\Pr[R:D=1]-\Pr[T:D=1]|
 \leq\min\{1,\epsilon_{\rm com}+\delta_{\rm nonce}\}.
\]
The corresponding $[-1,1]$ score bound is
$\min\{2,2(\epsilon_{\rm com}+\delta_{\rm nonce})\}$.
When the interface exposes no private-wire equality query, its additional
$\delta_{\rm nonce}$ term is zero.
\end{corollary}
\begin{proof}
First replace real byte equality by structural equality, coupling exactly
as above and charging $\delta_{\rm nonce}$. The resulting game has no extra
nonce-dependent equality observation, so apply the stated private-commitment
replacement there. Its ideal form is implemented exactly by the table
algorithm. A triangle inequality proves the bit bound. Applying that bound
to a test that returns one with probability $(1+Z)/2$ proves the score bound.
Without a private-wire equality operation the first replacement is absent.
\end{proof}

The hybrid order starts with constructible real record inputs, replaces
their private commitments, and then uses the table algorithm for the KEM
challenge embedding. Every changed commitment is consumed by that world's
subsequent intent prefixes, checks and outputs. For an adaptive origin-pair
event $E_j$, early stopping is justified only by a computable event making
$E_j$ false for every continuation and every still-admissible final target
choice attributed to $j$. The disjoint-event accounting above retains this
per-pair obligation.
